

\section{Basics}

\begin{frame}{Qu'est-ce que Cassandra\,?}

Cassandra, conçu à l'origine par Facebook, maintenant un projet de la fondation Apache,
distribué par la société Datastax.

\begin{redbullet}
  \item Initialement basé sur le système BigTable de Google (stockage orienté colonnes)
   \item A fortement évolué vers un \textit{modèle relationnel étendu}
  \item Un des rares systèmes NoSQL à proposer un typage fort
  \item Un langage de définition de schéma et de requêtes, CQL (pas de jointure)
    \end{redbullet}
   \vfill
   
Construit dès l'origine comme un système \alert{scalable} et \alert{distribué}.

\vfill

Ressource complémentaire\,: http://b3d.bdpedia.fr/cassandra.html

\end{frame}


\section{le modèle de données}


\begin{frame}
\frametitle{Tables et base de données}


Une \alert{base} est un ensemble de tables. On l'appelle \textit{keyspace}. Explication? Plus tard.

\vfill

Une \alert{table} est un ensemble de \textit{rows} conformes
au schéma de la table.

\vfill

Les tables sont \alert{dénormalisées}\,: on peut avoir des attributs dont
la valeur est un type construit (objet ou ensemble).
\vfill



Une ligne (\textit{row}) est un identifiant (\textit{row key}) associé à un ensemble de paires (clé, valeur).

\vfill

\alert{Perspective A}: c'est du relationnel étendu en rompant la première règle de normalisation (type atomique).

\vfill

\alert{Perspective B}: chaque ligne est un document structuré au sens où nous l'avons vu jusqu'ici.

\end{frame}



\section{Session découverte\,!}


\begin{frame}[fragile]
\frametitle{Application !}

Création d'un \textit{keyspace}.

\begin{minted}{sql}
create keyspace if not exists Pathologies 
       with replication = { 'class' : 'SimpleStrategy',
               'replication_factor': 3 };
\end{minted}

\vfill

Création d'une table des départements (sans imbrication).

\begin{minted}{sql}
 create table Departement (id text, 
                code text, 
                nom  int, 
                primary key (id) 
             );
\end{minted}
\end{frame}


\begin{frame}[fragile]
\frametitle{Insertions}

A la SQL

\begin{minted}{sql}
insert into Departement (id, code, nom)  
      values ('xa1', '16', 'Charentes');
insert into Departement (id,  code, nom)  
          values ('v6t', '42', 'Loire');
insert into Departement (id,  code, nom)  
          values ('9uj', '25', 'Doubs');
\end{minted}

\vfill

Avec JSON

\begin{minted}{sql}
insert into Departement JSON '{
				"id" : "uh6",
				"code": "43",
				"nom": "Haute-Loire"
			},' 
\end{minted}
\end{frame}



\begin{frame}[fragile]
\frametitle{Et les régions ?}

\vfill

On peut créer une autre table 

\begin{minted}{sql}
 create table Region (id text, 
                nom  int, 
                primary key (id) 
             );
\end{minted}

\vfil

\alert{mais on ne pourra pas faire la jointure avec le langage CQL...}. Il
faudrait écrire un programme...

\vfill

Deux solutions :
\begin{redbullet}
  \item Utiliser l'imbrication (données dénormalisées)
  \item Organiser physiquement les tables pour que le programme soit efficace
\end{redbullet}
\end{frame}


\begin{frame}[fragile]
\frametitle{Imbrication\,: la région dans les départements}

Création d'un type
\begin{minted}{sql}
create type region (id text,
                    nom text);
\end{minted}

\vfill

Table des départements, avec imbrication d'un objet conforme au type \texttt{region}.


\begin{minted}{sql}
 create table Departement (id text, 
                code text, 
                nom  int, 
                region frozen<region>,
                primary key (id) 
             );
\end{minted}

\vfill

On poura accéder aux régions, mais seulement en passant par les départements...
\end{frame}

\begin{frame}[fragile]
\frametitle{Insertion de documents}

À partir d'un document JSON

\begin{minted}{sql}
insert into Departement JSON '{
		"id" : "uh6",
		"code": "43",
		"nom": "Haute-Loire",
		"region": {
		     "id": "uin", 
		    "nom": "Auvergne-Rhône-Alpes"
		 }
	}' 
\end{minted}

\end{frame}


\begin{frame}[fragile]
\frametitle{Imbrication\,: les départements dans les régions}

Création d'un type pour le département
\begin{minted}{sql}
create type departement (id text,
                      "code": text,
                    nom text);
\end{minted}

\vfill

Table des régions avec imbrication \alert{d'une liste} d'objets conforme au type \texttt{departement}.


\begin{minted}{sql}
 create table Region (id text, 
                nom  int, 
                departements set<frozen<departement>>,
                primary key (id) 
             );
\end{minted}

\vfill

On poura accéder aux départements, mais seulement en passant par les régions...
\end{frame}

\begin{frame}[fragile]
\frametitle{Insertion de documents}

\begin{minted}{sql}
insert into Region JSON '{
		"_id": "3",
		"nom": "Auvergne-Rhônes-Alpes",
		"depts": [
			{"code": "01", "nom": "Ain"},
			{"code": "03", "nom": "Allier"}
			{"code": "42","nom": "Loire"}
		]
	}' 
\end{minted}

\end{frame}



\begin{frame}[fragile] \frametitle{Premiers pas avec CQL}

Sélectionnons tous les départements.

\begin{minted}{sql}
select  * from departement;
\end{minted}

\vfill

Ou seulement les 20 premiers.
\begin{minted}{sql}
select  * from departement limit 20;
\end{minted}

\vfill
Retour en JSON
\begin{minted}{sql}
select JSON * from departement;
\end{minted}

\vfill
Accédons à un objet imbriqué.
\begin{minted}{sql}
select nom, region.nom from departement;
\end{minted}

\vfill
En revanche, on ne peut pas accéder aux éléments d'un ensemble.
\begin{minted}{sql}
select nom, departements.nom from region;
\end{minted}
\end{frame}

\begin{frame}[fragile] \frametitle{La clause \texttt{where}}

Sélectionnons un film.

\begin{minted}{sql}
select  *  from region where id='iuv';
\end{minted}

Ou plusieurs

\begin{minted}{sql}
select  * from region 
    where id in ('iuv', 'r44', 'x45');
\end{minted}
 
Et si ce n'est pas la clé\,?

\begin{minted}{sql}
select  * from departement where nom='Loire' ;
\end{minted}

\begin{verbatim}
Unable to execute CQL script. Cannot execute this query 
as it might involve data
filtering and thus may have unpredictable performance. 
\end{verbatim}

Essayons de comprendre.
\end{frame}


%%%%%%%%%%
\begin{frame}[fragile] \frametitle{Restrictions sur les requêtes}

Toute recherche sur des critères autres qu'un préfixe de la clé entraîne un parcours
séquentiel

\vfill

Cassandra est conçu pour de très grandes bases de données, et le rejet 
de ces requêtes séquentielles est une précaution. 

\vfill

Principe: Cassandra limite les requêtes à celles dont le coût est proportionnel à la 
taille du résultat.

\vfill

On peut outrepasser avec la clause \texttt{allow filtering}.
\begin{minted}{sql}
select  * from departeme,nt 
     where nom='Loire' allow filtering;
\end{minted}


\end{frame}

\begin{frame}[fragile]
\frametitle{À retenir}

\begin{redbullet}
   \item Cassandra propose un modèle relationnel étendu, basé sur la capacité à imbriquer
     des types complexes dans la définition d'un schéma
   \item L'imbrication des constructeurs de type
   rend le modèle comparable aux documents structurés JSON ou XML.
   \item CQL est un pseudo-sql très limité : pas de jointure ou de capacité
       à interroger les structures imbriquées
   \item Le langage est conçu  pour accéder \alert{efficacement} à de très grosses
      bases, mais avec de très fortes limitations sur les requêtes
   \item Réciproquement, il faut concevoir le schéma pour tirer parti
   de ces principes\,: ce sera l'étape suivante
\end{redbullet}

\end{frame}